\section{Bài toán tối ưu}

\subsection{Mô hình baseline}

Mô hình baseline là checkpoint YOLO được huấn luyện theo cấu hình chuẩn trên tập
dữ liệu mục tiêu và suy luận ở FP32. Baseline cung cấp mốc chất lượng, kích thước
và hiệu năng để đánh giá mọi biến thể sau đó. Tên kiến trúc, phiên bản framework,
trọng số khởi tạo và cấu hình huấn luyện phải được cố định trong manifest thí
nghiệm.

\subsection{Hạn chế của mô hình baseline}

Baseline ưu tiên khả năng huấn luyện và chất lượng, chưa được tối ưu riêng cho
Jetson. Các hạn chế có thể gồm kích thước trọng số lớn, nhiều phép tính, latency
cao, GPU memory lớn và FPS chưa đạt yêu cầu. Mức độ hạn chế được xác nhận bằng
số đo trên cả Jetson Nano và Jetson Orin Nano thay vì suy luận chỉ từ FLOPs.

\subsection{Mục tiêu tối ưu}

Đồ án tối ưu đồng thời sáu tiêu chí: Accuracy, Latency, FPS, Model Size, GPU
Memory và Power. Bài toán được biểu diễn dưới dạng:
\begin{equation}
  \min_{m\in\mathcal{M}}
  \bigl(-A(m),\ L(m),\ -F(m),\ S(m),\ G(m),\ P(m)\bigr),
\end{equation}
trong đó $m$ là một cấu hình model, $A$ là accuracy, $L$ là latency, $F$ là FPS,
$S$ là kích thước, $G$ là GPU memory và $P$ là công suất. Thay vì gộp tùy ý
thành một điểm duy nhất, đồ án ưu tiên phân tích tập nghiệm Pareto và các ràng
buộc tối thiểu của ứng dụng.

Một cấu hình chỉ được xem là ứng viên triển khai khi mức giảm $mAP@50{:}95$
không vượt \cbs{ngưỡng chấp nhận}, latency đầu cuối không vượt \cbs{ngưỡng
thời gian thực} và thiết bị hoạt động ổn định trong bài đo kéo dài.